\section{Casos prácticos y gobernanza}

\subsection{Comparación de dos casos reales}

Considérense los procesos de reward learning de dos equipos representativos: OpenAI, en InstructGPT (2022), usó extensamente human preference data, mientras que el equipo de RL de Stanford, en tareas robóticas, se apoyó en goal classifiers. La siguiente tabla resume las diferencias entre ambos:

\begin{table}[H]
\centering
\begin{tabular}{p{0.25\textwidth} p{0.2\textwidth} p{0.2\textwidth} p{0.25\textwidth}}
\toprule
Dimensión & OpenAI (InstructGPT) & Stanford Robots & Conclusión \\
\midrule
Origen de los datos & preference comparison humana & success/failure snapshots & Las etiquetas humanas son más precisas, pero más costosas \\
Reward signal & sigmoid difference & binary classifier & Bradley-Terry es más adecuado para el diálogo; el goal classifier, para las deterministic tasks \\
Algoritmo de RL & PPO + KL & TRPO + KL flow & Ambos añaden una penalización de KL para evitar el drifting \\
Despliegue & API + guardrails & simulación → robot real & La API requiere una revisión de seguridad más amplia \\
\bottomrule
\end{tabular}
\caption{Comparación de dos procesos de reward learning}
\end{table}

\begin{warningbox}{Lecciones desde la perspectiva de la gobernanza}
Del lado de la API, OpenAI complementa el reward model con Rate limiting & guardrail filters; del lado de los robots, el equipo de Stanford depende más de instrumentation + simulators. Por lo tanto, la estrategia de gobernanza debe corresponder al tipo de tarea.
\end{warningbox}

\subsection{Observability Playbook}

Se recomienda firmemente construir el playbook de observabilidad según la siguiente checklist:

\begin{enumerate}
    \item Registrar para cada trajectory el reward model score + KL penalty
    \item Enviar las reward anomalies a un dashboard (por ejemplo, una sudden drop below 0.1)
    \item Detectar automáticamente la repetición excesiva, los fast reward jumps y los out-of-domain inputs
    \item Revisar cada semana los failure logs y consultar con el equipo de preferencias humanas si se necesitan más data
\end{enumerate}

\begin{importantbox}{Configuración de alertas automáticas}
Las reglas de alerta pueden basarse en la moving average del score del modelo: cuando la reward score window sube 3 std o tiene una sudden drop de 0.2 o más, se envía una notificación a PagerDuty y se hace trigger del rollback script.
\end{importantbox}

\subsection{Resumen de la sección}

Los casos prácticos muestran que el reward learning debe corresponder tanto a la tarea como a la gobernanza de la organización. La checklist y la observabilidad son esenciales para prevenir el reward hacking.
